% =============================================================================
%  The Harbor Volume — Workbook
%  Part V · The Anchor Protocol (Beam A — Identity & capability)
%  QUESTIONS
%
%  Self-contained \input fragment: the parent workbook supplies the document
%  class, the palette (hhcobalt/hhteal/hhamber/maydayred), and the \wbq /
%  \wbtier macros. The fallback definitions below let this file compile on its
%  own for proofing. Tiers: Check (recall) · Trace (derivation) · Open (design).
%  Difficulty: one to three filled stars.
% =============================================================================

\providecommand{\hhcobalt}{}  % no-op guard if palette not yet loaded
\providecolor{hhcobalt}{HTML}{003FB8}
\providecolor{hhteal}{HTML}{00564C}
\providecolor{hhamber}{HTML}{6B4500}
\providecolor{hhink}{HTML}{1B1712}

% Difficulty stars, set in small-caps cobalt so they read as a label, not ASCII.
\providecommand{\wbstars}[1]{{\footnotesize\textcolor{hhcobalt}{%
  \ifcase#1\or$\bigstar$\or$\bigstar\bigstar$\or$\bigstar\bigstar\bigstar$\fi}}}

% A tier tag: small-caps, weight 600, tracked — the eyebrow idiom, never clowny.
\providecommand{\wbtier}[1]{{\footnotesize\scshape\bfseries\textcolor{hhink}{#1}}}

% One question. #1 = local number, #2 = tier word, #3 = difficulty 1..3, #4 = body.
\providecommand{\wbq}[4]{%
  \par\medskip\noindent
  \makebox[0pt][r]{\textbf{V.#1}\hspace{1.2em}}%
  \wbtier{#2}~\wbstars{#3}\par\nopagebreak
  \noindent\ignorespaces #4\par}

\section*{Part V \quad The Anchor Protocol}
\addcontentsline{toc}{section}{Part V — The Anchor Protocol}

\noindent\textit{Beam A — Identity \& capability. Math-based identity for a
process: a verified control plane on localhost.} These problems exercise the
protocol's four phases (pinning, asymmetric identity, multi-hop delegation,
revocation), the three-layer verification stack, and the soundness boundary the
paper is careful to draw between the design proof and the deployed code.

% ───────────────────────────── Check tier ──────────────────────────────────
\wbq{1}{Check}{1}{%
State the one rule that makes a verifier immune to JWT algorithm confusion, and
explain in one sentence why reading the token's \texttt{alg} header --- even to
\emph{validate} it --- reintroduces the vulnerability.}

\wbq{2}{Check}{1}{%
A Harbor Card is delegated A~$\to$~B~$\to$~C. List every quantity that must
shrink or stay equal from one hop to the next, and name the one quantity that
must be \emph{fresh} (never reused) at each hop.}

\wbq{3}{Check}{1}{%
The paper claims the four phases are ``not four proofs.'' Which three phases does
ProVerif mechanize, and what is done instead for the fourth? Why does the fourth
phase not need a symbolic proof of the same kind?}

\wbq{4}{Check}{2}{%
Distinguish the two senses in which the Anchor Protocol is ``verified at unbounded
scale.'' One is delivered by ProVerif; one is explicitly future work. Name both
and say which quantity each ranges over.}

% ───────────────────────────── Trace tier ──────────────────────────────────
\wbq{5}{Trace}{2}{%
A cuckoo filter uses approximately $\log_2(1/\epsilon)+3$ bits per entry. Derive
the per-card cost at $\epsilon=10^{-3}$, then the total memory for $10^5$ active
revocations. Confirm the paper's ``$\approx 160$\,KB'' figure and state the one
assumption that makes it gossip-cheap.}

\wbq{6}{Trace}{2}{%
The freshness bound is stated two ways: ``$O(\log m)$ gossip rounds'' and
expected wall-clock $\Delta(1+\ln m)$. For $m=10$ daemons at $\Delta=30$\,s,
compute the expected convergence time and reconcile it with the figure's
``$\approx 2$\,minutes.'' Where does the gap between expected and worst case come
from?}

\wbq{7}{Trace}{3}{%
Reproduce the v3-vacuity argument. Given the reflexive-only order
\texttt{is\_subset(c1,c1)=true}, show that the no-escalation query holds for a
reason unrelated to the protocol. Then show precisely what
\texttt{harbor\_card\_v5} adds (the order, the insider adversary, the
non-vacuity query) so that the same query becomes a real defense. What does the
negative control establish?}

\wbq{8}{Trace}{3}{%
Construct the multi-hop escalation the single-hop proof misses:
$A{:}\texttt{write}\to B{:}\texttt{read}\to C{:}\texttt{write}$. Show why a
verifier that checks the \emph{final} capability against the \emph{root} accepts
it, and write the per-hop predicate that rejects it. Tie this to the runtime
requirement on $\mathrm{att}_B\circ\mathrm{att}_A$.}

% ───────────────────────────── Open tier ───────────────────────────────────
\wbq{9}{Open}{3}{%
The delegation-depth bound is a spawn-time policy check, not a verification-time
invariant. Sketch how you would extend the ProVerif model to unbounded chain
depth via recursive replication, and state the obstacle that makes this harder
than ProVerif's existing unbounded-\emph{session} guarantee.}

\wbq{10}{Open}{3}{%
The gossip-convergence bound currently leans on the standard anti-entropy
analysis rather than a mechanized proof, and the constant-time comparator's
side-channel resistance is inherited from an audited library. Pick one. Propose a
verification approach that would move it from \wbtier{open} to \wbtier{built},
and name the proof tool and the property you would discharge.}

\wbq{11}{Open}{2}{%
A Dolev--Yao adversary observing all gossip traffic learns which agents have been
disciplined. The paper offers salted-hash revocation IDs with daily salt rotation
as a mitigation. Analyse the residual leakage, and discuss what an adversary who
records gossip across many salt epochs can still infer.}
